\documentclass[a4paper,11pt]{article}

\usepackage[utf8]{inputenc}

\usepackage{minted}
\usepackage{graphicx}
\usepackage{placeins}

\begin{document}


\title{
    \textbf{Heap report}
}
\author{Jonathan Karlsson}
\date{Fall 2023}

\maketitle

\section*{Introduction}
This assignment consisted of first creating a prority queue as a linked list, then creating a heap as a tree, and finally turn the tree into an array 
which makes it easier to navigate.


\section*{Implementation}
\subsection*{Prority queue list}
The prority queue list was created as two, the first where the add method was simpler and the second where the remove method was simpler. On the first one the add method simply added
the value to the start of the queue, but the remove method is more complicated, as it goes through every element of the list, and then checks for the smallest element. When the samllest element is found it is removed in a similar way to a double linked list. Then the value is returned.
\begin{minted}{java}
public Integer remove() {
    if (this.first == null)
        return null;
    Node priority = this.first;
    Node next = this.first;
    while(next != null) {
        if(priority.head > next.head) {
            priority = next;
        }
    next = next.next;
    }
    if (priority.prev != null)
        priority.prev.next = priority.next;
    if (priority.prev == null)
        this.first = this.first.next;
    if(priority.next != null)
        priority.next.prev = priority.prev;

    return priority.head;
}    
\end{minted}
The second implementation, instead had a simple remove, where it only removed the first element, however here the add is more complicated as it has to go through the list until something larger is found and then the value is inserted infront of that larger value. 
This works by either finding the value in the middle where it gets inserted between two elements, or it is at the first element, where it gets inserted infront of the first, or finally it is not smaller than any value and then it is added to the end.
\begin{minted}{java}
public void add(Integer value) {
    if (this.first == null) {
        this.first = new Node(value, null, null);
        return;
    }
    Node next = this.first;
    Node prev = null;
    while (next != null) {
        if(next.head>value){
            Node nodeToAdd = new Node(value,next,next.prev);
            if (next.prev != null){
                next.prev.next = nodeToAdd;
                next.prev = nodeToAdd;
            }
            else {
                this.first = nodeToAdd;
            }
            break;
        }
        if(next.next == null) {
            prev = next;
        }
            next = next.next;
    }
    if (next == null) {
        prev.next = new Node(value, null, prev);
    }
}
\end{minted}
\subsection*{Tree heap}
The three heap was implemented with three methods add, remove and push. 
The add method works by first checking if the root is null, if it is then it adds the value to the root. If it was not null 
the add method calls another add method with the root, this add method recursivly swaps values downword until ther is a free space, it also works towards the nodes with the smallest size.
\begin{minted}{java}
private void add(Integer value,TreeNode node) {
    node.size +=1;
    if(value < node.value) {
        Integer temp = node.value;
        node.value = value;
        value = temp;
    }
    if (node.left == null) {
        node.left = new TreeNode(value);
        return;
    }
    if (node.right == null) {
        node.right = new TreeNode(value);
        return;
    }
    if(node.left.size < node.right.size) {
        add(value, node.left);
    }
    else {
        add(value,node.right);
    }
}
\end{minted}
The remove method starts by checking if there is no root, then null is returned, or root is alone, then root is set to null and its value is returned. If root was not null its value is saved, then
differnet if statements is run to check if the right or left value should be brought up, then current is set to the node which value was brought up and then the loop is run again until a node which does not have any nodes under it is reached and 
then it is set to null and the roots value is returned.
The push method only pushes adds a number to the root then works in a similar way to add and remove to swap the past root down until it is in a correct place. 
\begin{minted}{java}
public int push(Integer incr) {
    root.value += incr;
    TreeNode current = root;
    int depth = 0;
    while(((current.right!=null&&current.right.value < current.value)||
    (current.left!=null&& current.left.value < current.value))) {
        if(current.right!=null &&current.right.value<current.value) {
            Integer temp = current.right.value;
            current.right.value = current.value;
            current.value = temp;
            current = current.right;
            depth++;
        }
        else if(current.left!=null && current.left.value<current.value) {
            Integer temp = current.left.value;
            current.left.value = current.value;
            current.value = temp;
            current = current.left;
            depth++;
        }
    }
}
\end{minted}
\subsection*{ArrayHeap}
The add method for the arrayHeap is simple it just increases k and then adds the element to the k position and finally calls the swapInto method with k. 
The swapInto method works by first checking if root has been reached, if it has return, then checking if the parent to the currentIndex is -2/2 or -1/2. Finally if the current element is 
smaller than the parent it swaps and then calls the swapInto method with the parent. 
\begin{minted}{java}
private void swapInto(int currentIndex) {
    if(currentIndex <= 0)
        return;
    int parent;
    if(currentIndex%2 == 0) {
        parent = (currentIndex-2)/2;
    }
    else {
        parent = (currentIndex-1)/2;
    }
    if(array[currentIndex]<array[parent]) {
        Integer temp = array[currentIndex];
        array[currentIndex] = array[parent];
        array[parent] = temp;
        swapInto(parent);
    }
    
}
\end{minted}
The remove method works by first checking if the element in the root position is null if it is return null. Then the roots value is saved, the element in position k is put into root and the position k is set to null, then k is decreased and sink with root is called.
The sink method works by first calcuclating where the left and right index is. Then it checks if it can go into the left index, which it can if the leftindex is in the array, the leftindex is smaller than the current index. Then if the rightindex is not in the array or the rightindex element is larger than the leftindex elemnt. then the element in the current index and left index are swapped and sink is called with left index.
Then the same checks is done for the right index.
\begin{minted}{java}
private void sink(int currentIndex) {
    int leftIndex = currentIndex*2+1;
    int rightIndex = currentIndex*2+2;

    if(leftIndex <= k &&array[currentIndex]>array[leftIndex]) {
        if(rightIndex >k || array[leftIndex]<=array[rightIndex]){
            Integer temp = array[currentIndex];
            array[currentIndex] = array[leftIndex];
            array[leftIndex] = temp;
            sink(leftIndex);
        }
    }
    else if(rightIndex <= k && array[currentIndex]>array[rightIndex] ) {
        if(leftIndex > k || array[leftIndex]>array[rightIndex]) {
        Integer temp = array[currentIndex];
        array[currentIndex] = array[rightIndex];
        array[rightIndex] = temp;
        sink(rightIndex);  
        }
    }
}
\end{minted}

\section*{Benchmark}
The first benchmark was to compare the two linked lists. with each other and their add and remove functions. Here it seems that the first priority list with the complicated remove is slower, as the complicated add of the second list is 
much faster than the remove, while the simple add and simple remove is about as fast meaning that simple add + complicated remove is slower.
\newline
\begin{table}[h]
    \begin{center}
    \scalebox{1}{
    \begin{tabular}{l|c|c|c|c|}
    \textbf{n} & \textbf{L1 add({\textmu}s)} & \textbf{L2 add({\textmu}s)}& \textbf{L1 remove({\textmu}s)}& \textbf{L2 remove({\textmu}s)}\\
    \hline
    100    &   9  &     8&      14 &      3 \\
    200    &   7  &    19&      54 &      3 \\
    400    &   5  &    11&     215 &      2 \\
    800    &   9  &   141&     878 &      5 \\
   1600    &  14  &  1297&    3706 &     10 \\
    \end{tabular}
    }
    \caption{The result between the two linked lists combined}
    \label{tab:table1}
    \end{center}
\end{table}
\newline
\FloatBarrier{\noindent}
The second benchmark was to use the push operation a couple of times on the tree and compare it to remove and add with the same change of value. Then also save the amount of depth that the operation produced. 
The amount of times the push operation took was about 2,9{\textmu}s and the amount of time spent on the add and remove was 5,2{\textmu}s for 20 operations.
The avereage depth when pushing was around 3,35.
\newline
\newline
The final benchmark was to benchmark the array implementation and compare it to the linked implementation. As the second list is faster than the first list that 
is what is compared to the array. The result is clear that the array is much faster on a combined time, however the simple remove still beats the array, but as the array has a faster add it seems like the optimal choice.
\begin{table}[h]
    \begin{center}
    \scalebox{1}{
    \begin{tabular}{l|c|c|c|c|}
    \textbf{n} & \textbf{A add({\textmu}s)} & \textbf{L2 add({\textmu}s)}& \textbf{A remove({\textmu}s)}& \textbf{L2 remove({\textmu}s)}\\
    \hline
    100    &   9  &     8&      4 &      3 \\
    200    &   13  &    19&      7 &      3 \\
    400    &   23  &    11&     11 &      2 \\
    800    &   44  &   141&     23 &      5 \\
   1600    &  81  &  1297&    44 &     10 \\
    \end{tabular}
    }
    \caption{The result between the second list and the array}
    \label{tab:table2}
    \end{center}
\end{table}








\end{document}
